\documentclass[11pt]{article}
\usepackage[a4paper, margin=1in]{geometry}
\usepackage{parskip}
\usepackage{hyperref}
\hypersetup{colorlinks=true, urlcolor=blue}
\pagestyle{empty}

\begin{document}

\noindent
Radu-Tudor P\^irvu \\
Bucharest, Romania \\
+40 771 478 290 \\
radupirvu09@gmail.com

\vspace{1.5em}
\noindent
Ceragon Networks \\
Verification \& Validation Team \\
Calea Victoriei, Bucharest, Romania

\vspace{1.5em}
\noindent
Dear Hiring Team,

I am writing to apply for the Verification \& Validation Internship within Ceragon's R\&D department. I am a second-year student in Electronics, Telecommunications and Information Technology at Politehnica University of Bucharest, entering my third year this fall, after spending my 4th semester at Christian-Albrechts-Universit\"at zu Kiel. What draws me to this role specifically is the mindset your team is built around: trying to break a product before a customer ever gets the chance to. That's the way I like to work, and I would welcome the chance to apply it to actual wireless hardware instead of simulated coursework.

Two of my recent projects map closely onto what this role asks for. For my Biometric Parameter Monitor, I built an ESP32-based system that read BPM and SpO2 over I2C and transmitted the data to a remote server via HTTP, buffering and averaging sets of ten readings to filter sensor noise and reduce network overhead, a small-scale version of the kind of signal integrity and data reliability questions a V\&V engineer deals with daily. Separately, through coursework and personal projects I have used Wireshark to inspect and debug network traffic, which lines up directly with the networking fundamentals this role calls for. My background in embedded systems (Arduino, ESP32) and competitive robotics (FIRST Tech Challenge) has also left me comfortable with the build-test-debug cycle that testing and validation work runs on.

Outside the technical side, I have spent the last few years in structured, responsibility-heavy volunteer roles,x coordinating teams of up to twenty volunteers across five editions of The Color Run, and currently supporting adaptive climbing therapy sessions with Climb Again. Both have sharpened the kind of attention to detail and reliability under process that I imagine matters when documenting test results or tracking down unexpected behavior.

Thank you for considering my application. I look forward to the possibility of contributing to your team.

Sincerely, \\
Radu-Tudor P\^irvu

\end{document}
